% Options for packages loaded elsewhere
\PassOptionsToPackage{unicode}{hyperref}
\PassOptionsToPackage{hyphens}{url}
\documentclass[aspectratio=169,ignorenonframetext]{beamer}
\usepackage{ma2003b-beamer}
\hypersetup{
  pdftitle={Principal Component Analysis (PCA)},
  pdfauthor={Juliho Castillo Colmenares},
  hidelinks,
  pdfcreator={LaTeX via pandoc}}

\title{\texorpdfstring{Principal Component Analysis
(PCA)}{Principal Component Analysis (PCA)}}
\author{\texorpdfstring{Juliho Castillo
Colmenares}{Juliho Castillo Colmenares}}
\subtitle{Dimensionality Reduction, Spectral Decomposition, and Factor Modeling}
\institute{Tec de Monterrey}
\date{}

\begin{document}
\frame{\titlepage}

\begin{frame}{Today's Agenda}
\begin{enumerate}
\item
  Motivations \& Use Cases for Principal Component Analysis
\item
  Geometric Foundations: Rotation \& Variance Maximization
\item
  Spectral Decomposition of Covariance (\(\mathbf{S}\)) and Correlation
  (\(\mathbf{R}\))
\item
  Mathematical Derivation via Lagrange Multipliers
\item
  Component Retention Criteria (Kaiser, Scree, Parallel Analysis)
\item
  Loadings, Scores, and Communalities
\item
  High-Dimensional 2D/3D Biplots
\item
  Practical Case Study: Global Financial Asset Allocation
\end{enumerate}
\end{frame}

\begin{frame}{What is Principal Component Analysis?}
\textbf{Core Objective:} Transform \(p\) correlated variables
\(X_{1},\ldots,X_{p}\) into \(k \ll p\) uncorrelated variables
\(Y_{1},\ldots,Y_{k}\) (\textbf{Principal Components}) that capture the
maximum possible variance.

\textbf{Key Applications:}

\begin{itemize}
\item
  \textbf{Data Compression:} Compress feature space with minimal
  information loss.
\item
  \textbf{Multicollinearity Cure:} Replace correlated regressors with
  orthogonal components.
\item
  \textbf{Exploratory Visualization:} Project high dimensions onto 2D/3D
  score planes.
\item
  \textbf{Risk Modeling:} Identify macroeconomic drivers across assets.
\end{itemize}
\end{frame}

\begin{frame}{Geometry: Finding Axes of Maximum Variance}
{\def\LTcaptype{none} % do not increment counter
\begin{longtable}[]{@{}
  >{\raggedright\arraybackslash}p{(\linewidth - 2\tabcolsep) * \real{0.5455}}
  >{\raggedright\arraybackslash}p{(\linewidth - 2\tabcolsep) * \real{0.4545}}@{}}
\toprule\noalign{}
\endhead
\textbf{1. First Principal Component (\(Y_{1}\)):}
\(Y_{1} = \mathbf{a}_{1}^{T}\mathbf{X} = a_{11}X_{1} + \ldots + a_{1p}X_{p}\)
Direction of maximum data dispersion subject to
\(|\mathbf{a}_{1}| = 1\).

\textbf{2. Second Principal Component (\(Y_{2}\)):}
\(Y_{2} = \mathbf{a}_{2}^{T}\mathbf{X}\quad\text{ with  }\mathbf{a}_{2}^{T}\mathbf{a}_{1} = 0\)
Direction of maximum remaining variance, orthogonal to \(Y_{1}\). &
\textbf{Principal Axes of Ellipsoid:} Eigenvectors specify the
orientation of the constant-density ellipsoid axes. Half-lengths are
proportional to \(\pm c\sqrt{\lambda_{j}}\). \\
\bottomrule\noalign{}
\end{longtable}
}
\end{frame}

\begin{frame}{Mathematical Derivation (Lagrange Multipliers)}
Maximize
\(\operatorname{Var}(Y_{1}) = \mathbf{a}_{1}^{T}\mathbf{\Sigma}\mathbf{a}_{1}\)
subject to \(\mathbf{a}_{1}^{T}\mathbf{a}_{1} = 1\).

\textbf{Lagrangian Objective:}
\[\mathcal{L}(\mathbf{a}_{1},\lambda_{1}) = \mathbf{a}_{1}^{T}\mathbf{\Sigma}\mathbf{a}_{1} - \lambda_{1}\left( \mathbf{a}_{1}^{T}\mathbf{a}_{1} - 1 \right)\]

\textbf{First-Order Condition:}
\[\frac{\partial\mathcal{L}}{\partial\mathbf{a}_{1}} = 2\mathbf{\Sigma}\mathbf{a}_{1} - 2\lambda_{1}\mathbf{a}_{1} = \mathbf{0} \Longrightarrow \mathbf{\Sigma}\mathbf{a}_{1} = \lambda_{1}\mathbf{a}_{1}\]

\textbf{Result:} \(\lambda_{1}\) is the largest eigenvalue of
\(\mathbf{\Sigma}\), and \(\mathbf{a}_{1}\) is the corresponding
normalized eigenvector.
\end{frame}

\begin{frame}{Spectral Decomposition \& Variance Explained}
\textbf{Spectral Decomposition:}
\[\mathbf{S} = \mathbf{V}\mathbf{\Lambda}\mathbf{V}^{T} = \sum_{j = 1}^{p}\lambda_{j}\mathbf{v}_{j}\mathbf{v}_{j}^{T}\]

where
\(\mathbf{V} = \left\lbrack \mathbf{v}_{1},\ldots,\mathbf{v}_{p} \right\rbrack\)
is the eigenvector matrix and
\(\mathbf{\Lambda} = \operatorname{diag}(\lambda_{1},\ldots,\lambda_{p})\).

\textbf{Total Variance Preservation:}
\[\text{ Total Sample Variance } = \operatorname{tr}(\mathbf{S}) = \sum_{j = 1}^{p}s_{jj} = \sum_{j = 1}^{p}\lambda_{j}\]

\textbf{Proportion of Explained Variance:}
\[\text{ Prop}_{j} = \frac{\lambda_{j}}{\sum_{k = 1}^{p}\lambda_{k}}\]
\end{frame}

\begin{frame}{How Many Components to Retain?}
{\def\LTcaptype{none} % do not increment counter
\begin{longtable}[]{@{}
  >{\raggedright\arraybackslash}p{(\linewidth - 2\tabcolsep) * \real{0.5000}}
  >{\raggedright\arraybackslash}p{(\linewidth - 2\tabcolsep) * \real{0.5000}}@{}}
\toprule\noalign{}
\endhead
\begin{minipage}[t]{\linewidth}\raggedright
\textbf{1. Kaiser-Guttman Rule:}

\begin{itemize}
\item
  Retain \(\lambda_{j} > 1.0\) (for correlation \(\mathbf{R}\)).
\item
  Retains components explaining more variance than a single standardized
  variable.
\end{itemize}

\textbf{2. Cumulative Variance Threshold:}

\begin{itemize}
\item
  Retain components until cumulative variance reaches \(70\% - 85\%\).
\end{itemize}
\end{minipage} & \begin{minipage}[t]{\linewidth}\raggedright
\textbf{3. Cattell Scree Plot:}

\begin{itemize}
\item
  Visually identify the ``elbow'' where eigenvalue decrease flattens
  out.
\end{itemize}

\textbf{4. Horn's Parallel Analysis:}

\begin{itemize}
\item
  Retain components whose eigenvalues exceed the 95th percentile of
  simulated random noise.
\end{itemize}
\end{minipage} \\
\bottomrule\noalign{}
\end{longtable}
}
\end{frame}

\begin{frame}{Loadings, Scores, and the PCA Biplot}
\textbf{Component Loadings (Correlations with Variables):}
\[L_{jk} = \operatorname{Corr}(X_{k},Y_{j}) = a_{jk}\sqrt{\lambda_{j}}\quad\left( \text{for standardized data} \right)\]

\textbf{The Gabriel Biplot (1971):} Simultaneously displays:

\begin{itemize}
\item
  \textbf{Observation Points:} \(Y_{i1},Y_{i2}\) mapped in component
  score space.
\item
  \textbf{Variable Vectors:} Directed arrows showing loading magnitude
  and pairwise cosine correlation.
\end{itemize}
\end{frame}

\begin{frame}{Case Study: Multi-Asset Portfolio Decomposition}
\textbf{Setting:} 10 global asset classes and industry sector returns
(600 trading sessions).

\textbf{3 Retained Macroeconomic Factors (\(76.8\%\) Variance):}

\begin{enumerate}
\item
  \textbf{PC1 (Market Equity Beta - \(45.2\%\)):} Broad stock market
  risk (high loadings across US, Global, Tech).
\item
  \textbf{PC2 (Duration / Interest Rate - \(19.4\%\)):} Positive for
  Financials; negative for Treasuries \& REITs.
\item
  \textbf{PC3 (Commodity Inflation - \(12.2\%\)):} Driven by Gold \&
  Energy sector returns.
\end{enumerate}

\textbf{Takeaway:} Compress 10 risk dimensions into 3 uncorrelated macro
factors for hedging.
\end{frame}

\begin{frame}{Summary \& Key Takeaways}
\begin{itemize}
\item
  PCA is the cornerstone of unsupervised multivariate feature
  extraction.
\item
  Formulated via the \textbf{spectral decomposition} of \(\mathbf{S}\)
  or \(\mathbf{R}\).
\item
  Always standardize data (\(z\)-scores) when variables have different
  units or variances.
\item
  Validate component retention using multiple criteria (Kaiser, Scree,
  Parallel Analysis).
\item
  Use \textbf{Biplots} for holistic interpretation of observations and
  features.
\end{itemize}

{ \textbf{Next Chapter: Chapter 4 --- Factor Analysis (Latent Variable
Models)} }
\end{frame}

\end{document}




